The Linear Representation Hypothesis posits that high-level concepts are encoded as directions in the representation spaces of large language models. While this has been demonstrated for properties like truth, sentiment, and language identity, no prior work has investigated whether membership in a specific fictional franchise is similarly represented. We extract a ``Pok\'{e}mon direction'' from GPT-2-XL's residual stream using the mass-mean probing technique and find that the concept ``X is a Pok\'{e}mon'' is linearly encoded, achieving 88.4\% classification accuracy (AUC = 0.934) on held-out entities---far above the 50.5\% random baseline. By projecting 177 non-Pok\'{e}mon entities onto this direction, we discover a continuous spectrum of ``Pok\'{e}mon quality'': invented Pok\'{e}mon-sounding names, Nintendo game creatures (Cactuar, Koopa, Kirby), and Digimon score highest, while common real-world objects score lowest. The category-level differences are highly significant (Kruskal-Wallis $p = 1.03 \times 10^{-10}$). Thirteen non-Pok\'{e}mon entities cross the classification boundary, including six game creatures and one Digimon. These results extend the Linear Representation Hypothesis to fine-grained fictional franchise membership, revealing that language models encode not a binary ``is/isn't Pok\'{e}mon'' distinction but a graded, interpretable measure of Pok\'{e}mon-likeness.